% GENERATED FILE. Edit canonical lessons, then run npm run generate:books.
% canonical-source-hash: fnv1a64:7d06636b0a66a37a
% canonical-lessons: PT-W258-four-lines

\chapter{Writing --- Four Lines, And The Joiner That Points Back}
\label{ch:pt-four-lines}
\hlchaptermodality{258}

\begin{chapteropening}
\textbf{By the end of this chapter:} \emph{I can write four connected lines to a named reader, adding with e and giving a result with por isso, with the reason placed before it.}

Complete PT-W258-four-lines: four lines to one reader, at least one joined with e and one with por isso, then check the reason stands before por isso.
\end{chapteropening}

\section[por isso]{Four lines, and the joiner that points back}

\label{lesson:PT-W258-four-lines}



\begin{quote}
The last lesson asked for writing \textbf{finished on time}. This one asks for writing that \textbf{holds together}.

Pick one person who will read it.
\end{quote}

\subsection*{What to know first}
Connected writing is not \textbf{longer} writing. Four short lines can be connected; one long line can fail to be. What makes it connected is that a line stops standing on its own.

You hold more joiners than this lesson needs. Use two: \textbf{e} adds, and \textbf{por isso} gives a result. This book has not yet given you a word for \emph{because}; \textbf{por isso} does that job from the other side.

Notice which way it points. \textbf{por isso} is \textquotedblleft{}for that\textquotedblright{}, and the \emph{isso}, \textquotedblleft{}that\textquotedblright{}, points \textbf{back} at what came before it. So the reason comes \textbf{first}, and what follows \textbf{por isso} is what the reason led to.

\begin{quote}
\textbf{Moro em Lisboa, por isso falo português.} — \textquotedblleft{}I live in Lisbon, that is why I speak Portuguese.\textquotedblright{}
\end{quote}

Swap the two halves and it is still good Portuguese, but now speaking Portuguese is why you live in Lisbon. Nothing in the grammar warns you. Only the order carries the meaning.

\subsection*{Writing — connected composition}
Write \textbf{four lines} to the person you picked.

At least one must use \textbf{e}. At least one must use \textbf{por isso}, with the reason before it and the result after it.

Use only words you already write, and not the Lisbon line above. This lesson adds no vocabulary; the whole exercise is the joining.

\begin{tcolorbox}[breakable,colback=teal!4,colframe=teal!35!black,title={Before you move on}]
Cover your first line and read the second alone. Does it still make sense without the first? If it does, you wrote a \textbf{list}, which is where you were two lessons ago.

Then find your \textbf{por isso} line and say what led to what. If the reason sits after \textbf{por isso}, you have written the reverse of what you meant.
\end{tcolorbox}
